%% LyX 2.5.2 created this file.  For more info, see https://www.lyx.org/.
%% Do not edit unless you really know what you are doing.
\documentclass[12pt,spanish]{report}
\usepackage[T1]{fontenc}
\usepackage[utf8]{inputenc}
\setcounter{secnumdepth}{3}
\setcounter{tocdepth}{3}
\usepackage[modifiers.spanish=notilde.noquoting.nodecimaldot]{babel}
\usepackage{amsmath}
\usepackage{amssymb}
\usepackage[a4paper]{geometry}
\geometry{verbose}
\usepackage{setspace}
\onehalfspacing
\usepackage[bookmarks=false,
 breaklinks=false,pdfborder={0 0 1},backref=false,colorlinks=false]
 {hyperref}

\makeatletter
%%%%%%%%%%%%%%%%%%%%%%%%%%%%%% User specified LaTeX commands.


% ============================================================
% PAQUETES
% ============================================================
\PassOptionsToPackage{spanish,es-tabla}{babel}
\usepackage{babel}
\usepackage{graphicx}
\usepackage{booktabs}
\usepackage{array}
\usepackage{csquotes}

% ============================================================
% DATOS DE LA TESIS  (TODO: completar)
% ============================================================
\newcommand{\tesisTitulo}{Modelo de aprendizaje automático para la identificación y clasificación de patrones espaciales de criminalidad en el área urbana de Trujillo}
\newcommand{\tesisAutor}{Jefferson Miguel Peña Serrano y Martin Aryan Robles Perez}
\newcommand{\tesisAsesor}{Nombre del Asesor}
\newcommand{\tesisUniversidad}{Universidad Nacional de Trujillo}
\newcommand{\tesisFacultad}{Facultad de ...}
\newcommand{\tesisCarrera}{Escuela Profesional de ...}
\newcommand{\tesisTituloProfesional}{Ingeniero ...}
\newcommand{\tesisFecha}{2026}

% ============================================================
% DOCUMENTO
% ============================================================


\makeatother

\usepackage[style=apa,backend=biber,sorting=nyt]{biblatex}
\addbibresource{referencias/bibliografia.bib}
\begin{document}
% ============================================================
% PORTADA  (TODO: ajustar universidad, facultad, carrera, etc.)
% ============================================================
\begin{titlepage} \centering\par {\LARGE \tesisUniversidad}\\[0.3cm]
{\large \tesisFacultad}\\[0.3cm] {\large \tesisCarrera}\\[3.5cm]\par {\LARGE
\textbf{TESIS}}\\[1.2cm]\par {\LARGE\bfseries \tesisTitulo}\\[2.5cm]\par {\large
Para optar el título profesional de \tesisTituloProfesional}\\[2.5cm]\par {\large\textbf{Autor:}
\tesisAutor}\\[0.3cm] {\large\textbf{Asesor:} \tesisAsesor}\\[3.5cm]\par {\large
Trujillo -- Perú}\\[0.3cm] {\large \tesisFecha}\par \end{titlepage}


\tableofcontents{}\listoffigures
\listoftables

% ============================================================
% RESUMEN  (TODO: escribir el resumen final al terminar la tesis)
% ============================================================


\chapter*{Resumen}

\addcontentsline{toc}{chapter}{Resumen}

El presente trabajo tiene como objetivo desarrollar un modelo de aprendizaje
automático para la identificación y clasificación de patrones espaciales
de criminalidad en el área urbana de Trujillo. (TODO: completar con
objetivo, método, datos y resultados principales).

\vspace{1cm}

\noindent\textbf{Palabras clave:} aprendizaje automático, patrones
espaciales, criminalidad, análisis espacial, Trujillo.

\chapter*{Abstract}

\addcontentsline{toc}{chapter}{Abstract}

The aim of this thesis is to develop a machine learning model for
the identification and classification of spatial crime patterns in
the urban area of Trujillo. (TODO: translate the final abstract).

\vspace{1cm}

\noindent\textbf{Keywords:} machine learning, spatial patterns, crime,
spatial analysis, Trujillo.


% Cada carpeta = capítulo (numeración mayor: 1., 2., ...)
% Cada archivo dentro = sección (numeración menor: 1.1, 1.2, ...)

\chapter{Introducción}

\label{cap:introduccion}

El presente capítulo presenta el contexto, el problema, la justificación,
los objetivos y la hipótesis del estudio. (TODO: desarrollar este
párrafo introductorio).


\section{Antecedentes}

\label{sec:antecedentes}

% TODO: desarrollar esta sección con estudios previos
Los estudios sobre criminalidad urbana en América Latina han mostrado
que los delitos tienden a concentrarse en determinadas zonas de la
ciudad. Diversos autores han aplicado técnicas de análisis espacial
y aprendizaje automático para identificar estos patrones y apoyar
la toma de decisiones de las autoridades. (TODO: citar estudios previos
concretos).
 
\section{Planteamiento del problema}

\label{sec:planteamiento}

% TODO: desarrollar esta sección
La ciudad de Trujillo ha registrado en los últimos años un incremento
de la criminalidad, cuyas manifestaciones se distribuyen de forma
desigual en el territorio. Las autoridades no cuentan con herramientas
que permitan identificar y clasificar de manera automática los patrones
espaciales de los delitos. (TODO: formular la pregunta de investigación).


\section{Justificación}

\label{sec:justificacion}

% TODO: desarrollar esta sección
La presente investigación se justifica porque un modelo de aprendizaje
automático permitiría a las autoridades policiales y municipales asignar
mejor sus recursos, anticipar zonas de riesgo y diseñar estrategias
de prevención basadas en evidencia. (TODO: agregar justificaciones
teórica, práctica y social).
 \section{Objetivos}
\label{sec:objetivos}

% TODO: completar y ajustar los objetivos según el estudio

\subsection{Objetivo general}

Desarrollar un modelo de aprendizaje automático para la identificación y clasificación de patrones espaciales de criminalidad en el área urbana de Trujillo.

\subsection{Objetivos específicos}

\begin{itemize}
    \item Analizar la distribución espacial de los eventos delictivos en el área urbana de Trujillo.
    \item Seleccionar las variables espaciales y temporales relevantes para la detección de patrones de criminalidad.
    \item Entrenar y evaluar modelos de aprendizaje automático para la clasificación de los patrones identificados.
    \item Validar el modelo propuesto mediante métricas de rendimiento y comparación con enfoques tradicionales.
\end{itemize}
\section{Hipótesis}
\label{sec:hipotesis}

% TODO: completar según el estudio
Los algoritmos de aprendizaje automático aplicados a los datos georreferenciados de eventos delictivos permiten identificar y clasificar patrones espaciales de criminalidad en el área urbana de Trujillo con una precisión superior a los métodos de análisis espacial tradicionales. (TODO: ajustar la redacción de la hipótesis).
 \chapter{Marco Teórico}
\label{cap:marco_teorico}

En este capítulo se desarrollan los fundamentos teóricos sobre aprendizaje automático, análisis espacial del delito y el estado del arte de la clasificación de patrones de criminalidad. (TODO: desarrollar).

\section{Fundamentos teóricos}
\label{sec:fundamentos}

% TODO: desarrollar cada subsección con citas
\subsection{Aprendizaje automático}

El aprendizaje automático es una rama de la inteligencia artificial que permite a los modelos aprender patrones a partir de datos sin ser programados explícitamente \parencite{hastie2009elements}. Entre los algoritmos más utilizados para clasificación se encuentran los árboles de decisión, los bosques aleatorios y las máquinas de vectores de soporte \parencite{geron2019hands}.

\subsection{Análisis espacial del delito}

La criminología ambiental estudia la relación entre el delito y el espacio, señalando que los eventos delictivos no se distribuyen de manera uniforme sino que se concentran en puntos calientes o ``hot spots'' \parencite{eck2005mapping}. Estos patrones pueden identificarse mediante técnicas de estadística espacial y sistemas de información geográfica (SIG).

\subsection{Patrones espaciales de criminalidad}

Los patrones espaciales de criminalidad describen la forma en que los delitos se agrupan, se repiten o se desplazan en el territorio. Su identificación requiere transformar los datos de incidentes en variables espaciales como densidad, cercanía o vecindario, y luego aplicar métodos de clasificación \parencite{ratcliffe2008intelligence}.

\section{Estado del arte}

\label{sec:estado_del_arte}

% TODO: ampliar con trabajos recientes (2019-2026) sobre predicción de delitos con ML
Diversos estudios han aplicado modelos supervisados y no supervisados
para predecir y clasificar la criminalidad en entornos urbanos. Los
métodos de regresión y clasificación clásicos siguen siendo ampliamente
usados por su interpretabilidad \parencite{james2013introduction},
mientras que las redes neuronales profundas han mostrado buenos resultados
cuando se dispone de grandes volúmenes de datos georreferenciados
\parencite{goodfellow2016deep}.

En el contexto peruano existen pocas aplicaciones documentadas de
estos modelos a nivel de distritos urbanos, lo que constituye una
oportunidad para el presente trabajo. (TODO: revisar literatura nacional
y local).

\chapter{Metodología}
\label{cap:metodologia}

En este capítulo se describe el enfoque, el diseño, la población, la muestra y las técnicas e instrumentos utilizados en la investigación. (TODO: desarrollar).


\section{Enfoque y diseño de la investigación}

\label{sec:enfoque}

% TODO: desarrollar esta sección
La investigación tiene un enfoque cuantitativo y un diseño de tipo
descriptivo-explicativo, ya que busca describir la distribución espacial
de la criminalidad y explicar los patrones mediante un modelo de aprendizaje
automático. (TODO: justificar el enfoque y el diseño).


\section{Población y muestra}

\label{sec:poblacion}

% TODO: desarrollar esta sección
La población está conformada por los eventos delictivos registrados
en el área urbana de Trujillo durante el período de estudio. La muestra
estará constituida por los registros georreferenciados disponibles
de fuentes oficiales (policiales o municipales). (TODO: precisar cantidad
de registros y período).

\section{Técnicas e instrumentos de recolección de datos}
\label{sec:tecnicas}

% TODO: desarrollar esta sección
Para la recolección de datos se utilizarán fuentes secundarias de registros de incidentes delictivos con coordenadas geográficas. El procesamiento y modelamiento se realizará con el lenguaje Python y las bibliotecas de aprendizaje automático como scikit-learn, complementadas con herramientas de análisis espacial (SIG). (TODO: detallar instrumentos, software y procedimiento). 
\chapter{Resultados}

\label{cap:resultados}

En este capítulo se presentan los resultados obtenidos por el modelo
de aprendizaje automático. (TODO: desarrollar y colocar aquí las tablas
y figuras).

\section{Análisis de resultados}
\label{sec:analisis_resultados}

% TODO: desarrollar esta sección con los resultados reales
\begin{figure}[htbp]
    \centering
    % TODO: colocar la figura en la carpeta "figuras/" y descomentar:
    % \includegraphics[width=0.8\textwidth]{figuras/mapa_patrones.png}
    % \caption{Mapa de patrones espaciales de criminalidad en Trujillo.}
    % \label{fig:mapa_patrones}
\end{figure}

Los resultados del modelo muestran la identificación de los principales patrones espaciales de criminalidad. (TODO: describir métricas de rendimiento, precisión, matriz de confusión, etc.).

\chapter{Discusión}
\label{cap:discusion}

En este capítulo se contrastan los resultados obtenidos con los antecedentes y la literatura revisada. (TODO: desarrollar).


\section{Comparación con otros estudios}

\label{sec:comparacion}

% TODO: desarrollar esta sección
Los patrones espaciales identificados en Trujillo coinciden parcialmente
con los reportados en otras ciudades, donde los delitos se concentran
en zonas de alta actividad comercial y de transporte. (TODO: comparar
métricas y métodos con otros trabajos, y discutir limitaciones).
 \chapter{Conclusiones}
\label{cap:conclusiones}

En este capítulo se presentan las conclusiones y recomendaciones del estudio. (TODO: desarrollar).

\section{Conclusiones y recomendaciones}
\label{sec:conclusiones}

% TODO: desarrollar esta sección
Se concluye que el modelo de aprendizaje automático permitió identificar y clasificar los patrones espaciales de criminalidad en el área urbana de Trujillo. (TODO: redactar cada conclusión ligada a los objetivos específicos y añadir recomendaciones para autoridades y trabajos futuros).

\printbibliography[heading=bibintoc,title=Referencias bibliográficas]

\end{document}
